\section{Related Works}
\label{sec:related}

\textbf{Structural Representations for Manipulation.}
Structural representations determine the orchestration of different modules in a manipulation system and yield different implications on the capabilities, assumptions, efficiency, and effectiveness of the system.
Rigid-body poses are most commonly used given well-understood rigid-body motions in free space and their efficiency at modeling long-range dependencies of objects~\cite{kaelbling2010hierarchical,kaelbling2013integrated,srivastava2014combined,byravan2017se3,dantam2018incremental,migimatsu2020object,garrett2021integrated,curtis2022long,labbe2022megapose,tyree20226,pan2023tax,wen2023foundationpose}.
However, since it often requires both geometry and dynamics of environment to be modeled beforehand, various works have studied structural representations using data-driven methods, such as learning object-centric representation~\cite{lenz2015deepmpc,chang2016compositional,battaglia2016interaction,sanchez2018graph,jang2018grasp2vec,tremblay2018deep,xu2019densephysnet,mao2019neuro,burgess2019monet,hewing2020learning,locatello2020object,heravi2023visuomotor,zhu2023learning,yuan2022sornet,cheng2023nod,hsu2024s}, particle-based dynamics~\cite{li2018learning,lin2022planning,wang2023dynamic,shi2023robocook,lin2022learning,abou2024physically,bauer2024doughnet}, and keypoints or descriptors~\cite{schmidt2016self,florence2018dense,manuelli2019kpam,kulkarni2019unsupervised,qin2020keto,sundaresan2020learning,manuelli2020keypoints,chen2021unsupervised,simeonov2022neural,simeonov2023se,vecerik2023robotap,chun2023local,bahl2023affordances,wen2023any,bharadhwaj2024track2act}.
Among them, keypoints have shown great promises given their interpretability, efficiency, generalization to instance variations~\cite{manuelli2019kpam}, and ability to model both rigid bodies and deformable objects.
However, manual annotation is required per task, thus lacking scalability in open-world settings, which we aim to address in this work.






\textbf{Constrained Optimization in Manipulation.}
Constraints are often used to impose desired behaviors on robots. Motion planning algorithms use geometric constraints to compute feasible trajectories that avoid obstacles and achieve goals~\cite{kingston2018sampling,ratliff2009chomp,schulman2014motion,sundaralingam2023curobo,marcucci2024shortest,ratliff2018riemannian}.
Contact constraints can be used to plan forceful or contact-rich behaviors~\cite{posa2016optimization,mordatch2012discovery,mordatch2012contact,posa2014direct,howell2022predictive,liu2024model,lynch1996stable,hou2018fast,sleiman2023versatile,yang2024dynamic,graesdal2024towards,yunt2006trajectory,yunt2007combined,yunt2011augmented}.
For sequential manipulation tasks, task and motion planning (TAMP) \cite{kaelbling2010hierarchical,kaelbling2013integrated,garrett2021integrated} is a widely used framework, often formulated as constraint satisfaction problems~\cite{lagriffoul2012constraint,lagriffoul2014efficiently,lozano2014constraint,dantam2018incremental,yang2023compositional,silver2021learning,vu2024coast} with continuous geometric problems as subroutines.
Logic-Geometric Programming \cite{toussaint2015logic,toussaint2017multi,toussaint2018differentiable,ha2020probabilistic,xue2023d,driess2020deep} alternatively formulates a nonlinear constrained program over the entire state trajectory, taking into account both logical and geometric constraints.
Constraints can either be manually written or learned from data in the form of manifolds~\cite{sutanto2021learning}, feasibility model \cite{driess2020deep,yang2022sequence}, or signed-distance fields~\cite{driess2022learning,camps2022learning}.
Inspired by \cite{toussaint2022sequence}, we formulate sequential manipulation tasks as an integrated \emph{continuous} mathematical program that is repeatedly solved in a receding-horizon fashion, with the key difference being that the constraints are synthesized by foundation models.



\textbf{Foundation Models for Robotics. }
Leveraging foundation models for robotics is an active area of research.
We refer readers to~\cite{hu2023toward,firoozi2023foundation,kawaharazuka2024real,yang2023foundation} for overview and recent applications.
Here, we focus on VLMs that are capable of incorporating visual inputs~\cite{radford2021learning,ramesh2021zero, li2022blip,achiam2023gpt,li2023blip,openai2023gpt} for manipulation.
However, despite showing promises for open-world planning and goal specification~\cite{huang2024copa,liu2024moka,nasiriany2024pivot,hu2023look,du2023video,hong20233d,chen2024spatialvlm,huang2023voxposer,brohan2023rt,gao2023physically,wang2024grounding,hsu2023ns3d,gao2024physically,yuan2024robopoint,duan2024manipulate}, the caption-guided pretraining scheme of VLMs often limits visual details that can be retained about the image~\cite{tong2024eyes,thrush2022winoground,yuksekgonul2023and,hsieh2024sugarcrepe}.
Self-supervised vision models (e.g., DINO~\cite{caron2021emerging, oquab2023dinov2}), on the other hand, provide fine-grained pixel-level features useful for various vision and robotic tasks~\cite{amir2021deep,melas2022deep,wang2023d3fields,zhu2023learning,lin2023spawnnet,di2024keypoint,di2024dinobot}, but there lack effective ways for interpreting open-world semantics pivotal for cross-task generalization.
In this work, we leverage their complementary strengths by using DINOv2~\cite{oquab2023dinov2} for fine-grained keypoint proposal and using GPT-4o~\cite{openai2023gpt} for its visual reasoning capability in a supported output modality (code). %
Similar forms of visual prompting techniques are also explored in concurrent works~\cite{huang2024copa,liu2024moka,nasiriany2024pivot,yuan2024robopoint,lee2024affordance}.
In this work, we demonstrate~\algabrvname possesses the unique advantages of performing challenging 6-12 DoF tasks, integrated high-level reasoning for reactive replanning, high-frequency closed-loop execution, and generating black-box constraints via visual prompting. More discussion in Appendix~\ref{sec:concurrent}.










